\documentclass[11pt,a4paper]{article}
\usepackage[T1]{fontenc}
\usepackage{isabelle,isabellesym}
\usepackage{amsmath,amssymb,url}

% This should be the last package used.
\usepackage{pdfsetup}

\urlstyle{rm}
\isabellestyle{it}

\begin{document}

\title{The Isoperimetric Inequality for Rectifiable Simple Closed Curves}
\author{Lawrence C. Paulson}
\date{}
\maketitle

\begin{abstract}
This entry formalises the planar isoperimetric inequality for rectifiable
simple closed curves. If a curve has length $L$, the area of its inside is
at most $L^2/(4\pi)$; equality implies that its image is a circle. The
development proves a Wirtinger inequality, a special area formula for
convex curves, and a convexification argument. It builds on the
\texttt{Fourier} AFP entry and analysis material in Isabelle2026.
\end{abstract}

\tableofcontents

\section{Proof outline and provenance}

The proof first establishes the inequality for convex curves. Arc-length
reparametrisation and the Wirtinger inequality supply its analytic core;
a special case of Green's area formula relates the resulting integral to
the enclosed area. Convexification then reduces a general rectifiable
simple closed curve to the convex case and accounts for equality.

This development was ported from John Harrison's HOL Light formalisation
\texttt{100/isoperimetric.ml}.\footnote{\url{https://github.com/jrh13/hol-light/blob/master/100/isoperimetric.ml}}
The arc-length material also follows
\texttt{Multivariate/integration.ml}.\footnote{\url{https://github.com/jrh13/hol-light/blob/master/Multivariate/integration.ml}}
Claude assisted with the translation; the Isabelle proofs required
substantial manual revision. There were no additional textual sources.
The Wirtinger proof uses the AFP entry
\emph{Fourier Series}.\footnote{\url{https://isa-afp.org/entries/Fourier.html}}

\input{session}

\end{document}
